%----------SUMMARY----------%
\section{Summary}
\begin{itemize}[leftmargin=0.15in, label={}]
    \small{\item{
    I design, build and operate an MLOps platform on on-prem Kubernetes. Over six months at HUINNO I codified the entire path — ingestion $\rightarrow$ lakehouse $\rightarrow$ transformation $\rightarrow$ lineage $\rightarrow$ distributed training $\rightarrow$ experiment / model registry — in \textbf{Pulumi}, as the \textbf{sole owner of cluster construction} on a two-person platform team (93\% of stack-code changes and 83\% of the ETL repository's commits are mine). I remove human bottlenecks rather than add tools, and decide by measurement. Earlier: iOS $\rightarrow$ backend $\rightarrow$ frontend $\rightarrow$ cloud data pipelines, then led a data engineering team.
    }}
\end{itemize}

%----------KEY RESULTS----------%
\section{Key Results (measured 2026-09)}
\begin{itemize}[leftmargin=0.15in, label={}]
    \small{
    \item{\textbf{Built the entire MLOps pipeline} — ingestion through the model registry, with no gap left manual, codified into 12 Pulumi stacks: infrastructure changes land as PRs and a new environment is rebuilt by composing stacks. (1,204 commits, 678 merged PRs.)}
    \item{\textbf{10 nodes / 235 running pods / 23 internal services} in production. SSO across all 23 with Dex (GitHub OIDC) + oauth2-proxy, the same identity carried into Kubernetes RBAC.}
    \item{\textbf{Removed the single-person dependency in dataset preparation} — "ask one data scientist and wait" became scheduled Airflow 3.x + dbt + Apache Iceberg runs with lineage in DataHub; researchers now query directly. (1,599 commits in the ETL repo; I wrote the shared framework myself, along with a product pipeline, the MIMIC-IV pipeline and the operational ones.)}
    \item{\textbf{AI agents let one person carry what normally takes a team} — I wrote most of this Kubernetes platform with coding agents running continuously, and have run it in production for six months. That speed is held honest by 719 SDK unit tests, CI gates, 23 incident write-ups and 30 PRDs.}
    \item{\textbf{Earlier:} Led the AWS EMR / PySpark art-market ETL pipeline (Eazel, team lead).}
    }
\end{itemize}

%----------CURRENT WORK----------%
\section{What I Do Now}
\begin{itemize}[leftmargin=0.15in, label={}]
    \small{\item{
    Data Engineer, HUINNO (Dec 2025 --). I operate every layer of the MLOps platform on bare-metal Kubernetes and am attaching a reproducibility contract (version, content hash, code commit, environment lock and seed stamped onto a dataset at run time) plus MLflow tracking and a model registry, so research code can move onto KubeRay. A quantified survey found the root cause — results having no coordinates — and I am establishing the split where researchers own domain ETL and the platform owns execution, lineage, verification and retries.
    }}
\end{itemize}

%----------STACK----------%
\section{Stack}
\begin{itemize}[leftmargin=0.15in, label={}]
    \small{\item{
    \textbf{Platform}: Kubernetes (kubeadm) · Pulumi · Cilium · SeaweedFS · Dex/OIDC · Prometheus/Grafana \\[2pt]
    \textbf{Data}: Airflow 3.x · dbt · Apache Iceberg · DataHub · PostgreSQL (pg\_duckdb) · ClickHouse \\[2pt]
    \textbf{ML}: Ray/KubeRay · MLflow · PyTorch · DVC \\[2pt]
    \textbf{Languages \& other}: Python · SQL · Swift · Shell · AWS/GCP · LangChain
    }}
\end{itemize}

%----------CERTIFICATIONS----------%
\section{Certifications \& Languages}
\begin{itemize}[leftmargin=0.15in, label={}]
    \small{\item{
    \textbf{Certifications}: Engineer Information Processing \mbox{(Nov 2017)} · Craftsman Information Processing \mbox{(Jul 2009)} — issued by HRD Korea \\[2pt]
    \textbf{Languages}: Korean (native) · English \mbox{OPIc AL} · Japanese (conversational)
    }}
\end{itemize}

%----------DETAIL----------%
% One sentence does not warrant a section heading; keep it as a footer line
\vspace{6pt}
{\small Full experience and project details are in a separate document (\texttt{resume\_english.pdf}).}
